\addsec*{Abstract}


Engineering work is strongly shaped by structured processes that require accuracy, traceability, and compliance with standards. Typical tasks such as handling change requests, performing risk analyzes, or preparing test plans follow detailed steps, need background information at the right time, and end with proper documentation. In practice, these activities are often time-consuming, spread across different tools, and result in documentation that varies in quality. With recent progress in large language models (LLMs), there is new potential to support engineers by giving process guidance, helping with information retrieval, and partly automating reporting. At the same time, the use of such models in industry is still in an early stage and faces open questions regarding process awareness, reliability, and organizational control.

This thesis looks at how LLMs can be applied in engineering processes in a way that improves both efficiency and quality while still meeting the requirements of organizations. The work will begin by looking at challenges in industry, using existing process documents as input. The goal is to identify pain points and set clear improvement goals. After that, the current state of research on AI process support will be reviewed, with particular attention to methods like retrieval-augmented generation (RAG) and structured prompting strategies.

The findings of this first phase will be used to build a domain model that captures the main aspects of process support, such as activities, roles, artifacts, constraints, and decision points. This model will serve as the basis for a prototype prototype of the proof-of-concept that demonstrates AI-assisted, process-aware assistance. The prototype will be designed to give engineers step-by-step guidance, provide relevant knowledge when needed, and create documentation of process execution. The evaluation will compare the prototype with current practices, looking at efficiency, correctness, compliance, and acceptance by users.

The expected contributions of this work are: (1) a taxonomy of industrial challenges in process support, (2) a domain model for LLM-assisted process guidance, (3) a proof-of-concept system, and (4) an evaluation of its potential benefits and limitations.




\textbf{\textit{Keywords: Large Language Models(LLMs), Retrieval-Augmented Generation (RAG), Process Guidance}} 
